\clearpage
\fancyhead[L]{\small TAVOLE DELLE PARALLASSI.}
\hypertarget{native-604-lower}{}
\begin{center}\small Göttingische gelehrte Anzeigen. Stück 108. Seiten 1073..1077. 7. Juli 1817.\\[3pt]\rule{.72\linewidth}{.3pt}\end{center}

\emph{Tavole delle parallassi di altezza, di longitudine e di latitudine calcolate dagli astronomi dell' osservatorio dell' università Gregoriana nel collegio romano.} Rom 1816. Nella stamperia de Romanis. XXXIV und 122 Seiten im Folioformat.

\begingroup\brokenpenalty=10000\relax
Die Bedeckungen von Fixsternen durch den Mond und die Sonnenfinsternisse sind für die Astronomie und Geografie von so großer Bedeutung, dass jeder Beitrag zur Erleichterung der darauf bezogenen Berechnungen dankbar aufgenommen zu werden verdient. Die sogenannte Methode des Neunzigsten ist bei den Astronomen zur Berechnung dieser Phänomene noch immer am gebräuchlichsten, und die Berechnung der Längen- und Breitenparallaxe des Mondes macht einen Hauptteil dieser Methode aus. Die vorliegenden Tafeln sollen diese Berechnung der Parallaxen abkürzen. Der Plan dazu stammt von \textsc{Andreas Conti}, einem der Astronomen der römischen Sternwarte. Bei der Ausführung dieser umfangreichen Arbeit unterstützten ihn seine Kollegen \textsc{Joseph Calandrelli} und \textsc{Jacob Reichenbach}.

Damit sich unsere Leser eine Vorstellung vom Aufbau und Gebrauch dieser Tafeln und vom dadurch zu erzielenden Gewinn machen können, müssen wir sie einzeln genauer untersuchen. Es sind eigentlich \emph{drei} Tafeln. Die erste umfasst 93 Seiten, die zweite und die dritte je 12 Seiten. Dazu kommen auf einem eigenen Blatt noch einige kleine Hilfstafeln. Die erste Tafel enthält nur die Produkte aus der halben Horizontalparallaxe und den Kosinuswerten aller Winkel. Sie hat also zwei Argumente. Das eine Argument ist die Horizontalparallaxe, deren 50 Werte von $53'\,20''$ bis $61'\,30''$ in Schritten von 10 Sekunden fortlaufen. Das zweite Argument bilden alle Winkel des Quadranten in Schritten von 10 Minuten. Jede Seite ist zweigeteilt. Die größere Hälfte enthält die eben genannten Produkte, die andere, kleinere Hälfte gibt unter dem Titel Supplement der ersten Tafel die Produkte aus denselben Kosinuswerten und den halben Horizontalparallaxen der Planeten (von $1''$ bis $40''$). Außerdem gibt sie nochmals die Produkte aus diesen Kosinuswerten und jeder einzelnen Minute von $1'$ bis $9'$ sowie jedem Vielfachen von zehn Minuten von $10'$ bis $90'$. Wie man sieht, wurde also an Platz nicht gespart.

\SourceLeaf{605}Die zweite Tafel hat ebenfalls zwei Argumente. Sie gibt uns den Wert des Produkts $\alpha(\sec L - 1)$ für die Argumente $L$ und $\alpha$. Das erste durchläuft die Werte von 0 bis $6^{\circ}$ in Schritten von zehn Minuten, das zweite die Werte von 0 bis $60'\,30''$ in Halbminutenschritten und zusätzlich alle ungeraden Sekunden von $1''$ bis $35''$.

Die dritte Tafel enthält 12 Fixsterne, die vom Mond häufig bedeckt werden, den Wert eines Hilfswinkel $\varphi$, der teils von der Breite des Fixsterns, teils von einem weiter unten zu erklärenden Winkel $D - \tfrac{1}{2}\Pi$ abhängt. Die Werte dieses Winkels durchlaufen den ganzen Quadranten von 20 zu 2c Minuten.

Die drei kleinen Hilfstafeln enthalten außerdem den Unterschied zwischen der Tangente und dem Bogen, die Umrechnung von Sekunden in Dezimalteile der Minute und den Wert des Produkts, $\operatorname{tang} \tfrac{1}{2}(D+\tfrac{1}{2}\Pi)\cdot\sin L'$, wobei das erste Argument $D+\tfrac{1}{2}\Pi$ alle Winkel des Quadranten und das zweite, $L'$, die Winkel von 0 bis $6^{\circ}$ in Schritten von einem Grad umfasst.

Das ist der Inhalt der Tafeln. Uns bleibt noch zu erklären, wie sie gebraucht werden sollen.

Die Höhenparallaxe erhält man sofort, indem man das Ergebnis der ersten Tafel verdoppelt, wenn man die Horizontalparallaxe und die scheinbare Höhe als Argumente verwendet. Ist stattdessen die wahre Höhe gegeben, so wird man die gesuchte Größe durch wiederholte Annäherung erhalten. Dieselbe Bemerkung gilt auch für das Folgende, sofern die zu verwendenden Argumente nicht unmittelbar gegeben sind.

Die Längenparallaxe $\Pi$ wird stets hinreichend genau durch die Formel
\[
\Pi = \dfrac{P \sin h \sin(D+\Pi)}{\cos L}
\]
bestimmt, wobei $P$ die Horizontalparallaxe, $h$ die Höhe des Neunzigsten, $D$ die wahre Länge abzüglich der Länge des Neunzigsten und $L$ die wahre Breite bezeichnet. Ist $\alpha$ die Differenz der Ergebnisse der ersten Tafel, wenn man sie zunächst mit $P$ und $D-h+\Pi$ und dann mit $P$ und $D+h+\Pi$ benutzt, so wird $\Pi$ gleich der Summe sein aus $\alpha$ und dem Ergebnis der zweiten Tafel, wenn man diese mit den Argumenten $\alpha$ und $L$ benutzt.

Die Breitenparallaxe $\pi$ ergibt sich hinreichend genau aus der Formel
\[
\pi = -P[\cos h \cos(L+\pi) - \sin h \sin(L+\pi)\cos(D+\tfrac{1}{2}\Pi)]
\]
oder aus dieser
\[
\pi = -\dfrac{P\cos(h+\varphi)\cos(L+\pi)}{\cos\varphi}
\]
wobei $\varphi$ einen durch die Formel
\[
\operatorname{tang} \varphi = \operatorname{tang}(L+\pi)\cdot\cos(D+\tfrac{1}{2}\Pi)
\]
bestimmten Hilfswinkel bezeichnet. Da die letzte Formel für $\pi$ dieselbe Form hat wie die Formel für $\Pi$, sieht man leicht, wie sich die erste und die zweite Tafel zu ihrer Berechnung verwenden lassen.

Die Hilfsmittel $\varphi$ gibt die dritte Tafel, wenn die Bedeckung eines ihrer zwölf Sterne nach der Methode von \textsc{Carlini} berechnet werden soll. Allgemein aber findet man ihn, indem man das Supplement der ersten Tafel mit der ersten Hilfstafel verbindet. Da der Hilfswinkel $\varphi$ von der scheinbaren Breite abhängt, ist das folgende Verfahren vorgeschrieben, wenn stattdessen die wahre Breite gegeben ist. Setzt man wie oben
\[
P \sin h \sin(D+\Pi) = \alpha
\]
und zugleich
\SourceLeaf{606}\[
\sin(L+\pi)\,\operatorname{tang}\tfrac{1}{2}(D+\tfrac{1}{2}\Pi) = M
\]
so gilt hinreichend genau
\[
\pi = -P\cos(h+\lambda+\pi) - \alpha M
\]
wobei man $M$ mithilfe der letzten Hilfstafel und $\alpha$ und $P\cos(h+\lambda+\pi)$ mithilfe der ersten Tafel findet. Auch hier ist wiederholte Annäherung unvermeidlich.

Wenn wir nun unser Urteil über diese Tafeln offen aussprechen sollen, müssen wir gestehen, dass wir darin keine Erleichterung der Parallaxenrechnungen finden können. Wir halten das Rechnen mit den gewöhnlichen trigonometrischen Tafeln für mindestens ebenso bequem wie die Verwendung dieser Spezialtafeln. Druck und Papier sind übrigens schön, wie man es sich bei allen Zahlenwerken wünschen möchte.

Bemerkenswert ist auch die Widmung dieses Werks an den Papst. Man freut sich über die Unterstützung, die er der römischen Sternwarte gewährt. Was soll man aber von dem leidenschaftlichen Eifern der römischen Astronomen gegen „die Gottlosen, die die Astronomie auf die schändlichste und schrecklichste Weise missbrauchen, das Herz verderben, die wahren Gläubigen verhöhnen und sie als unsinnige Anbeter der Sonne verspotten“ halten? Wer versteht diese seltsamen Beschuldigungen, und wer hätte ohne einen unzweideutigen Hinweis in dieser Widmung erraten, dass sie einem großen Geometer gelten sollen, der gewagt hat, eine Hypothese zur physikalischen Erklärung der Bildung des Sonnensystems aufzustellen?
\par\clearpage\endgroup
